\documentclass[compacto,10pt]{pubciencias-notas}

% -- Paquetes adicionales
\usepackage{comandosEPN}

% -- Datos del libro
\autor{Departamento de Formación Básica}
\materia{Álgebra Lineal}
\nota{Semana 04 Clase 01}
\fecha{Semestre 2025-A}

\logouno[1.1cm]{Logo_EPN.eps}
\logodos[1.7cm]{Logo_DFB.eps}

\definecolor{colordef}{cmyk}{0.81,0.62,0.00,0.22}

\begin{document}

\encabezado

\vspace*{-8mm}

\setcounter{section}{1}
\section{Sistemas de ecuaciones}
\setcounter{subsection}{1}
\subsection{Resolución de sistemas lineales}

\begin{defi}
    Sean $A\in \Mat[\K]{m}{n}$ y $b\in \Mat[\K]{m}{1}$, dado el sistema de ecuaciones lineales
    \[
        Ax=b,
    \]
    se dice que
    \begin{itemize}
    \item
        el sistema es \textbf{inconsistente} si no tiene solución;
    \item
        el sistema es \textbf{consistente} si tiene al menos una solución.
    \end{itemize}
\end{defi}

\begin{ejem}\hspace{0pt} 
El sistema de ecuaciones lineales
\[
    \systeme[x_1 x_2]{x_1 -x_2 = 3,
        2x_1 -2x_2 = 4}
\]
no tiene solución, es decir, es un sistema inconsistente. Analicemos la matriz ampliada del sistema
\begin{align*}
    \begin{pmatrix}
    1&-1&|&3\\
    2&-2&|&4
    \end{pmatrix}
    &\sim
    \begin{pmatrix}
    1&-1&|&3\\
    0&0&|&-2
    \end{pmatrix} 
    && -2F_{1} + F_{2} \to F_{2}
\end{align*}
Notemos que la segunda fila de la matriz equivalente implica que
\[
    0x_{1} + 0x_{2} = -2
\]
lo cual es falso, lo que nos permite concluir que el sistema no posee solución, es decir, es inconsistente para cualquier valor de $x_{1}$ y $x_{2}$.
\end{ejem}

\begin{ejem}\hspace{0pt} 
El sistema de ecuaciones lineales
\[
    \systeme[x_1 x_2]{x_1 -x_2 = 3,
        2x_1 -2x_2 = 6}
\]
tiene al menos una solución, es decir, es un sistema consistente. Analicemos la matriz ampliada del sistema
\begin{align*}
    \begin{pmatrix}
    1&-1&|&3\\
    2&-2&|&6
    \end{pmatrix}
    &\sim
    \begin{pmatrix}
    1&-1&|&3\\
    0&0&|&0
    \end{pmatrix} 
    && -2F_{1} + F_{2} \to F_{2}
\end{align*}
Notemos que la segunda fila de la matriz equivalente implica que
\[
    0x_{1} + 0x_{2} = 0
\]
lo cual es verdadero para cualquier valor de $x_{1}$ y $x_{2}$, mientras que la primera fila implica que
\[
    x_{1} - x_{2} = 3
\]
que se verifica para todos los valores de $x_{1} = 3 + x_{2}$ con $x_{2}\in\R$. Por tanto, el sistema es consistente.
\end{ejem}

\begin{ejer}\hspace{0pt} 
Dado el sistema de ecuaciones lineales
\[
    \systeme[x_1 x_2]{2x_1 -x_2 = 3,
        x_1 +x_2 = 0.}
\]
Muestre que es un sistema consistente.
\end{ejer}

\begin{comment}
\begin{ejer}\hspace{0pt} 
Tipos de soluciones y aplicación del teorema de Rouché-Frobenius.
\begin{enumerate}
\item[$a)$] Considere el sistema de ecuaciones lineales
\[
    \systeme[x_1 x_2 x_3]{x_1 +x_2 +x_3 = 1,
        2x_1 +3x_2 = 1,
        x_1 +x_2 -x_3 = 0}
\]
Compruebe que se cumple la igualdad $\rang(A) = \rang(A|b) = 3$ y concluya que el sistema es consistente. ¿Cuántas soluciones posee el sistema? 
\item[$b)$] Dado el sistema de ecuaciones lineales
\[
    \systeme[x_1 x_2 x_3]{x_2 +x_3 = 1,
        2x_1 +3x_2 -x_3 = 1,
        x_1 +x_2 -x_3 = 0.}
\]
Demuestre que este sistema posee infinitas soluciones.
\end{enumerate}
\end{ejer}

\begin{ejer}\hspace{0pt} 
Dado $\alpha\in\R$, considere al sistema de ecuaciones lineales
\[
    \systeme[x_1 x_2]{{\alpha}x_1 +x_2 = 1,
        x_1 +{\alpha}x_2 = 1.}
\]
\begin{enumerate}
\item[$a)$] ¿Para qué valores de $\alpha$ el sistema tiene solución única?
\item[$b)$] ¿Para qué valores de $\alpha$ el sistema no tiene solución?
\item[$c)$] ¿Para qué valores de $\alpha$ el sistema tiene infinitas soluciones?
\end{enumerate}
\end{ejer}
\end{comment}

\begin{teo}
    Sean $A\in \Mat[\K]{m}{n}$ y $b\in \Mat[\K]{m}{1}$, dado el sistema de ecuaciones lineales
    \[
        Ax=b,
    \]
    se tiene una y solo una de las siguientes
    \begin{itemize}
    \item
        el sistema es inconsistente;
    \item
        el sistema es consistente y tiene una única solución; o
    \item
        el sistema es consistente y tiene infinitas soluciones.
    \end{itemize}
\end{teo}

\begin{teo}[Teorema de Rouché–Frobenius]\label{Sem03TeoR-F}
    Sean $A\in \Mat[\K]{m}{n}$ y $b\in \Mat[\K]{m}{1}$, dado el sistema de ecuaciones lineales
    \[
        Ax=b,
    \]
    se tiene que
    \begin{itemize}
    \item
        el sistema es consistente si y solo si $\rang(A)=\rang(A|b)$;
    \item
        en caso de que el sistema sea consistente, la solución es única si y solo si $\rang(A)=n$.
    \end{itemize}
\end{teo}

\begin{ejem}\hspace{0pt} 
Considere el sistema de ecuaciones lineales
\[
    \systeme[x_1 x_2 x_3]{-x_1 +x_2 +x_3 = 1,
        2x_1 +3x_2 -2x_3 = 1,
        x_1 +x_2 -x_3 = 0}
\]
es un sistema inconsistente. En efecto, se tiene que 
\[
    \rang(A) = 2 \neq \rang(A|b) = 3.
\]
\end{ejem}

\begin{obs} Notemos que podemos determinar el tipo de solución, lo cual lo hacemos analizando el rango de las matrices de coeficientes y ampliada, sin embargo, no hemos determinado la solución como tal.
\end{obs}
\end{document}